\chapter*{Einf"uhrung, Terminologie}
\addcontentsline{toc}{chapter}{\numberline{}{Einf"uhrung, Terminologie}}

\section*{Terminologie}
\addcontentsline{toc}{section}{\numberline{}{Terminologie}}

Das kleinst m"ogliche Modell in der Kryptographie hat folgende Form:

\begin{itemize}
 \item Sender Alice, Empf"anger Bob;
 \item Nachricht oder Klartext $m$, Chiffrat $c$;
 \item Verschl"usselung $E$ (encryption), Entschl"usselung $D$ (decryption).
\end{itemize}

Alice will Bob eine Nachricht $m$ schicken. Alice will erreichen,
da\3 ihre Nachricht sicher "ubertragen wird, d.h. ein Unbefugter
darf die Nachricht nicht lesen oder ver"andern k"onnen. 
Um dies zu erreichen, mu\3 die Nachricht von
Alice verschl"usselt werden, so da\3 sie nur von Bob
entschl"usselt werden kann.

\begin{center}
  \begin{tabular}{ccc}
              &   \shortstack{verschl"usselt\\ $E$ encryption}  &               \\
    Klartext  &   $\longrightarrow$                             & Chiffrat      \\
    (message) &                                                 & (ciphertext)  \\
    $m$       &   $\longleftarrow$                              & $c$           \\
              &   \shortstack{entschl"usselt\\ $D$ decryption}  & 
  \end{tabular}
\end{center}

Mathematisch kann der Sachverhalt wie folgt beschrieben
werden. Die Funktion $E$ ist eine injektive Funktion von der
Menge der Klartexte in die Menge der Chiffrate, $D$ ist eine
surjektive Funktion in der entgegengesetzten Richtung.

\begin{eqnarray*}
  E(m) & = & c \\
  D(c) & = & m \\
  D(E(m)) & = & m \\
  E(D(c)) & = & c 
\end{eqnarray*}


\section*{Aufgaben der Kryptographie}
\addcontentsline{toc}{section}{\numberline{}{Aufgaben der Kryptographie}}

{\em Kryptographie} oder auch {\em Kryptologie} (aus dem Griechischen
kryp\'os "`versteckt"' und l\'ogos "`Wort"') ist die Wissenschaft, die
sich mit der sicheren (geheimen) Kommunikation befa\3t. Die
Vertraulichkeit ist nur eine der Aufgaben, die die Kryptographie heute
erf"ullen mu\3. "Ublicherweise werden die folgenden vier Eigenschaften
betrachtet:
\begin{itemize}
  \item{\bf Vertraulichkeit:} \\Geheimhaltung der Nachricht.
  \item{\bf Integrit"at:}  \\ Empf"anger kann verifizieren, da\3 die Nachricht nicht ver"andert wurde.
  \item{\bf Verbindlichkeit:} \\ Unabstreitbarkeit, Sender kann die Herkunft nicht ableugnen.
  \item{\bf Authentizit"at:} \\ eindeutige Herkunft, Nachricht ist wirklich vom Sender.
\end{itemize}


%% *************** Algorithmen und Schl"ussel
\section*{Algorithmen und Schl"ussel}
\addcontentsline{toc}{section}{\numberline{}{Algorithmen und Schl"ussel}}

Ein {\em kryptographischer Algorithmus} ist die Realisierung einer
mathematischen Funktion, die zur Ver- und Entschl"usselnung benutzt
wird. (Im allgemeinen gibt es zwei verschiedene Funktionen zur Ver-
und Entschl"usselung.)


\subsection*{Algorithmen ohne Schl"ussel}
\addcontentsline{toc}{subsection}{\numberline{}{Algorithmen ohne Schl"ussel}}

\begin{itemize}
  \item Sicherheit beruht auf der Geheimhaltung des Algorithmus.
  \item Der Algorithmus kann schlecht auf G"ute "uberpr"uft werden, 
        falls kein kompetentes Mitglied in der Gruppe vorhanden ist, 
        das den Algorithmus benutzt.
  \item Personen in mehreren Gruppen ben"otigen mehrere Algorithmen.
  \item Nach einem erfolgreichen Angriff oder nachdem ein Mitglied die Gruppe
        verlassen hat, mu\3 der Algorithmus vollst"andig ge"andert werden.
\end{itemize}


\subsection*{Algoritmen mit Schl"ussel}
\addcontentsline{toc}{subsection}{\numberline{}{Algorithmen mit Schl"ussel}}

Es wird ein Algorithmus verwendet, der von Kryptographen gepr"uft oder
standardisiert werden kann. Die Sicherheit wird durch die Wahl eines
geheimen Schl"ussels erreicht.  Es gibt zwei Klassen von Algorithmen
mit Schl"ussel.
\begin{enumerate}
  \item{\bf Symmetrische Algorithmen}\\
    Die Schl"ussel zur Ver- und Entschl"usselung sind einfach auseinander berechenbar, 
    meistens sind sie identisch. Sender und Empf"anger einigen sich vor dem Senden 
    von Nachrichten auf den (die) Schl"ussel.
  \begin{itemize}
    \item {\bf  Stromchiffren:} Nachrichten werden bit- bzw. byteweise verarbeitet. 
          Das Chiffrat eines Zeichens h"angt von den vorigen Zeichen bzw. von der Stelle im
          Text ab.
    \item {\bf Blockchiffren:} Nachrichten  werden blockweise (64 bit, 128 bit) verarbeitet. 
          Ein Klartextblock wird unabh"angig von der Stelle seines Auftretens auf ein
          Chiffrat abgebildet.
  \end{itemize}
 \item{\bf Asymmetrische Algorithmen} (oder Public-Key-Algorithmen) \\
  Es werden unterschiedliche Schl"ussel verwendet. Der Schl"ussel zur Entschl"usselung 
  l"a\3t sich nicht mit vertretbarem Aufwand aus dem Schl"ussel zum
  Verschl"usseln berechnen. Der Schl"ussel zum Verschl"usseln hei\3t 
  {\em "offentlicher Schl"ussel}, der andere {\em geheimer Schl"ussel}.
\end{enumerate}


%% ************** Kryptoanalyse
\section*{Kryptoanalyse}
\addcontentsline{toc}{section}{\numberline{}{Kryptoanalyse}}

Das Ziel der {\em Kryptoanalyse} ist es, aus einem Chiffrat $c$ die zugeh"orige Nachricht $m$ 
ohne Kenntnis des geheimen Schl"ussels zu gewinnen oder den geheimen
Schl"ussel zu finden. Die Grundannahme ist dabei, 
da"s der Algorithmus bekannt ist. Es gibt vier Typen von Angriffen:
\begin{enumerate}
  \item{\bf Ciphertext-Only} \\
   \begin{tabular}{rl}
     Geg:  & Einige Chiffrate, die mit demselben Schl"ussel und Algorithmus \\
           & verschl"usselt wurden. \\
     Ziel: & Klartexte und/oder Schl"ussel zu finden.
   \end{tabular}
  \item{\bf Known-plaintext} \\
   \begin{tabular}{rl}
     Geg:  & Paare von Klartext und Chiffrat. \\
     Ziel: & Finden des Schl"ussels (um sp"ater unbekannte Klartexte zu finden).\\
   \end{tabular}
  \item{\bf Chosen-plaintext} \\
   \begin{tabular}{rl}
     Geg:  & Klartexte k"onnen frei gew"ahlt werden, deren Chiffrate dann berechnet\\
           & werden. \\
     Ziel: & Schl"ussel zu finden.
   \end{tabular}
  \item{\bf Chosen-ciphertext} \\
   \begin{tabular}{rl}
     Geg:  & Chiffrate k"onnen frei gew"ahlt werden, die dann entschl"usselt werden\\ 
           & (bei Public-Keyverfahren sinnvoll). \\
     Ziel: & Schl"ussel (geheimen Schl"ussel bei PK) zu finden.
   \end{tabular}
\end{enumerate} 
Alle Kryptosysteme k"onnen durch einen Known-plaintext-Angriff gebrochen werden, 
indem man alle m"oglichen Schl"ussel ausprobiert. Diese Art, den Schl"ussel aus einem
Klartext-Chiffrat-Paar zu gewinnen, hei\3t {\em Brute-force-Angriff}. Oft ist dieser 
Angriff wegen der Redundanz in der Nachricht (Sprache, bekannter Anfang, Anrede, usw. ) auch
m"oglich, wenn nur Chiffrate gegeben sind. Die Kryptographie sucht Algorithmen, die mit den zur 
Verf"ugung stehenden Mitteln (heutige oder zuk"unftige) nicht gebrochen werden k"onnen.

\section*{Sicherheit von Algorithemen}
\addcontentsline{toc}{section}{\numberline{}{Sicherheit von Algorithemen}}

Es gibt vier Kategorien, die das Versagen eines kryptographischen Algorithmus beschreiben.
\begin{enumerate}
\item{\bf Vollst"andiger Zusammenbruch} \\ 
  Der (geheime) Schl"ussel wurde gefunden.
\item{\bf Globale Entschl"usselung} (global deduction) \\
 Es wurde ein anderer Algorithmus zum Entschl"usseln gefunden, ohne den Schl"ussel zu kennen.
\item{\bf Lokale Entschl"usselung} \\
  Es wurden von einigen Chiffraten die Klartexte gefunden.
\item{\bf Informationsgewinn} \\
  Teile oder Information "uber den Klartext und/oder den Schl"ussel wurden gefunden.
\end{enumerate}


\section*{Aufwand von Angriffen}
\addcontentsline{toc}{section}{\numberline{}{Aufwand von Angriffen}}

Der Aufwand der Angriffe kann auf drei verschiedene Weisen gemessen werden:
\begin{enumerate}
  \item{\bf Rechenkapazit"at:} \\
    Zeit bzw. Anzahl der Rechenschritte.
  \item{\bf Speicherkomplexit"at:} \\
    ben"otigter Speicher.
  \item{\bf Datenkomplexit"at:} \\
    ben"otigte Anzahl an Chiffraten bzw. Chiffrat-Klartext-Paaren.
\end{enumerate}    
  
  
